\chapter{Physical Implementations} \label{app:physical_implementations}

\minitoc

\section{Physical Implementation of NOT Circuit}

\subsection{CMOS switches}

A logic gates consists in a network of interconnected switches
implemented using the two type of MOS transistors: p-MOS and
n-MOS. How behaves the two type of transistors in specific
configurations is presented in Figure \ref{0_switch}.

A switch connected to $V_{DD}$ is implemented using a p-MOS
transistor. It is represented in Figure \ref{0_switch}a {\em off}
(generating $z$, which means {\em Hi-Z}: no signal, neither 0, nor
1) and in Figure \ref{0_switch}b it is represented {\em on}
(generating 1 logic, or {\em truth}).

A switch connected to {\em ground} is implemented using a n-MOS
transistor. It is represented in Figure \ref{0_switch}c {\em off}
(generating $z$, which means {\em Hi-Z}: no signal, neither 0, nor
1) and in Figure \ref{0_switch}e it is represented {\em on}
(generating 0 logic, or {\em false}).

\placedrawing[h]{figures/02_02_transistors.lp}{{\bf Basic switches.} {\small {\bf a}.
Open switch connected to $V_{DD}$. {\bf b}. Closed switch connected to
$V_{DD}$. {\bf c}. Open switch connected to {\em ground}. {\bf d}.
Closed switch connected to {\em ground}.}}{0_switch}

A MOS transistor works very well as an {\em on-off} switch
connecting its drain to a certain potential.  A p-MOS transistor
can be used to connect its drain to a high potential when its
gates is connected to ground, and an n-MOS transistor can connect
its drain to ground if its gates is connected to a high potential.
This complementary behavior is used to build the elementary logic
circuits.

In Figure \ref{src} is presented the {\em switch-resistor-capacitor
model} (SRC). If $V_{GS} < V_{T}$ then the transistor is {\bf off}, if
$V_{GS} \geq V_{T}$ then the transistor is {\bf on}. In both cases the
input of the transistor behaves like a capacitor, the gate-source
capacitor $C_{GS}$.

When the transistor is on its drain-source resistance is:
$$R_{ON} = R_{n}\frac{L}{W}$$
where: $L$ is the channel length, $W$ is the channel width, and $R_{n}$
is the resistance per square. The length $L$ is a constant
characterizing a certain technology. For example, if $L = 0.13\mu m$
this means it is about a $0.13\mu m$ process.

The input capacitor has the value:
$$C_{GS} = \frac{\epsilon_{OX} LW}{d}.$$
The value:
$$C_{OX} = \frac{\epsilon_{OX}}{d}$$
where: $\epsilon_{OX} \approx 3.9 \epsilon_{0}$ is the permittivity of
the silicon dioxide, is  the gate-to-channel capacitance per unit area
of the MOSFET gate.

In this conditions the gate input current is:
$$i_{G} = C_{GS} \frac{dv_{GS}}{dt}$$

\placedrawing[h]{figures/02_03_mosfet.lp}{{\bf The MOSFET switch.} {\small The {\em
switch-resistor-capacitor} model consists in the two states: OF
($V_{GS} < V_{T}$), and ON ($V_{GS} \geq V_{T}$). In both states the
input is defined by the capacitor $C_{GS}$.}}{src}


\begin{exemplu}
For an AND gate with low strength, with $W = 1.8\mu m$, in $0.13
\mu m$ technology, supposing $C_{OX} = 4fF/\mu m^{2}$, results the
input capacitance: $$C_{GS} = 4 \times 0.13 \times 1.8 fF = 0.936
fF$$ Assuming $R_{n} = 5K\Omega$, results for the same gate:
$$R_{ON} = 5 \times \frac{0.13}{1.8} \; K\Omega = 361 \Omega$$

$\diamond$
\end{exemplu}

\subsection{The Inverter}


\subsubsection{The static behavior}

The smallest and simplest logic circuit -- the invertor -- can be built
using a pair of complementary transistors, connecting together the two
gates as input and the two drains as output, while the n-MOS source is
connected to ground (interpreted as logic 0) and the p-MOS source to
$V_{DD}$ (interpreted as logic 1). Results the circuit represented in
Figure \ref{0_inv}.

\placedrawing[h]{figures/02_04_inverter.lp}{{\bf Building an invertor.} {\small {\bf
a}. The invertor circuit. {\bf b}. The logic symbol for the invertor
circuit.}}{0_inv}

The behavior of the invertor consist in combining the behaviors of the
two switches previously defined. For $in=0$ pMOS is $on$ and nMOS is
$off$ the output generating $V_{DD}$ which means 1. For $in=1$ pMOS is
$off$ and nMOS is $on$ the output generating 0.

The static behavior of the inverter (or NOT) circuit can be easy
explained starting from the switches described in Figure
\ref{0_switch}. Connecting together a switch generating $z$ with a
switch generating 1 or 0, the connection point will generate 0 or
1.

\subsubsection{Dynamic behavior}

The propagation time of an inverter can be analyzed using the two
serially connected invertors represented in Figure \ref{switch}. The
delay of the first invertor is generated by its capacitive load,
$C_{L}$, composed by:
\begin{itemize}
\item its parasitic drain/bulk capacitance, $C_{DB}$, is the intrinsic  output capacitance of the first invertor
\item wiring capacitance, $C_{wire}$, which depends on the length of the wire (of width $W_w$ and of length $L_w$) connected between the two invertors: $$C_{wire} = C_{thickox} W_w L_w$$
\item next stage input capacitance, $C_{G}$, approximated by summing the gate capacitance for pMOC and nMOS transistors: $$C_G = C_{Gp} + C_{Gn} = C_{ox}(W_p L_p + W_n L_n)$$
\end{itemize}

\placedrawing[h]{figures/02_05_propagation_time.lp}{{\bf The propagation time.} {\small
}}{switch}

The total load capacitance
$$C_L = C_{DB} + C_{wire} + C_G$$
is sometimes dominated by $C_{wire}$. For short connections $C_G$ dominates, while for big {\em fan-out}\index{fan-out: the number on input gates  connected to the output of a digital circuit} both, $C_{wire}$ and $C_G$ must be considered.

The signal $V_A$ is used to measure the propagation time of the first NOT in Figure \ref{switch}a. It is generated by an ideal pulse generator with output impedance 0. Thus, the rising time and the falling time of this signal are considered 0 (the input capacitance of the NOT circuit is charged  or discharged in no time).

The two delay times (see Figure \ref{switch}c) associated to an invertor
(to a gate in the general case) are defined as follows:
\begin{itemize}
    \item $t_{pLH}$: the time interval between the moment the input
    switches in 0 and the output reaches $V_{OH}/2$ coming from 0
    \item $t_{pHL}$: the time interval between the moment the input
    switches in 1 and the output reaches $V_{OH}/2$ coming from $V_{OH}$
\end{itemize}

Let us consider the transition of $V_A$ from 0 to $V_{OH}$ at $t_r$ (rise edge). Before transition, at $t_r^-$, $C_L$ is fully charged and $V_B = V_{OH}$. In Figure \ref{switch}b is represented the equivalent circuit at $t_r^+$, when pMOS is off and nMOS is on. In this moment starts the process of discharging  the capacitance $C_L$ at the constant current
$$I_{Dn(sat)} = \frac{1}{2} \mu_n C_{ox} \frac{W_n}{L_n} (V_{OH} - V_{Tn})^2$$
In Figure \ref{switch1}, at $t_r^-$ the transistor is cut, $I_{Dn} = 0$. At $t_r^+$ the  nMOS transistor switch in saturation and becomes an ideal {\em constant current generator} which starts to discharge $C_L$ linearly at the constant current $I_{Dn(sat)}$. The process continue until $V_{OUT} = V_{OH}$, according to the definition of $t_{pHL}$.


\placedrawing[h]{figures/02_06_nmos_output.lp}{{\bf The output characteristic of the nMOS transistor.} {\small
}}{switch1}

In order to compute $t_{pHL}$ we take into consideration the constant value of the discharging current which provide a linear variation of $v_{OUT}$.
$$\frac{dv_{out}}{dt} = \frac{d}{dt}(\frac{q_L}{C_L}) = \frac{-I_{Dn(sat)}}{C_L}$$
$$\frac{dv_{out}}{dt} = \frac{\frac{V_{OH}}{2} - V_{OH}}{t_{pHL}}$$
We solve the equations for $t_{pHL}$:
$$t_{pHL} = C_L \frac{1}{\mu_n C_{ox} \frac{W_n}{L_n} (V_{OH} - V_{Tn})} \frac{V_{OH}}{V_{OH} - V_{Tn}}$$
Because:
$$R_{ONn} = \frac{1}{\mu_n C_{ox} \frac{W_n}{L_n} (V_{OH} - V_{Tn})}$$
results:
$$t_{pHL} = C_L R_{ONn} \frac{1}{1 - \frac{V_{Tn}}{V_{OH}}} = k_n R_{ONn} C_L = k_n \tau_{nL} $$
where:
\begin{itemize}
\item $\tau_{nL}$ is the {\em constant time} associated to the H-L transition
\item $k_n$ is a constant associated to the technology we use; it goes down when $V_{OH}$ increases or $V_T$ decreases
\end{itemize}

The speed of a gate depends by its dimension and by the capacitive load
it drives. For a big $W$ the value of $R_{ON}$ is small charging or
discharging $C_{L}$ faster.

For $t_{pLH}$ the approach is similar. Results: $t_{pLH} = k_p \tau_{pL}$.

By definition the propagation time associated to a circuit is: $$t_p = (t_{pLH} + t_{pHL})/2$$
its value being dominated by the value of $C_L$ and the size (width) of the two transistors, $W_n$ and $W_p$.



\section{Physical Implementation of NAND \& NOR Gates}

The 2-input AND circuit, $a \cdot b$, works like a ``gate" opened by the signal $a$ for the signal $b$. Indeed, the gate is ``open" for $b$ only if $a = 1$. This is the reason for which the AND circuit was baptised {\bf {\em gate}}\index{gate}. Then, the use imposed this alias as the generic name for any logic circuit. Thus, AND, OR, XOR, NAND, ... are all called {\em gates}.


\subsection{The static behavior of gates}


For 2-input NAND and 2-input NOR gates the same principle will be
applied, interconnecting 2 pairs of complementary transistors to
obtain the needed behaviors.

There are two kind of interconnecting rules for the same type of
transistors, p-MOS or n-MOS. They can be interconnected serially
or parallel.

A serial connection will establish an {\em on} configuration only
if both transistors of the same type are {\em on}, and the
connection is {\em off} if at least one transistor is {\em off}.

A parallel connection will establish an {\em on} configuration if
at least one is {\em on}, and the connection is {\em off} only if
both are {\em off}.

\placedrawing[h]{figures/A1_01_nand_gate.lp}{{\bf The  NAND gate.} {\small {\bf a}.
The internal structure of a NAND gate: the output is 1 when at least one input is 0.  {\bf b}. The logic symbol
for NAND.} }{0_nand}

Applying the previous rules result the circuits presented in
Figures \ref{0_nand} and \ref{0_nor}.



\placedrawing[h]{figures/A1_02_nor_gate.lp}{{\bf The NOR gate.} {\small {\bf a}. The
internal structure of a NOR gate: the output is 1 only when both inputs are 0. {\bf b}. The logic symbol for
NOR.} }{0_nor}

For the NAND gate the output is 0 if both n-MOS transistors are
{\em on}, and the output is one when at least on p-MOS transistor
is {\em on}. Indeed, if $A=B=1$ both $n$ transistors are {\em on}
and both $p$ transistors are {\em off}. The output corresponds
with the definition, it is 0. If $A=0$ or $B=0$ the output is 1,
because at least one $p$ transistor is {\em on} and at least one
$n$ transistor is {\em off}.

A similar explanation works for the NOR gate. The main idea is to
design a gate so as to avoid the simultaneous connection of
$V_{DD}$ and ground potential to the output of the gate.

For designing an AND or an OR gate we will use an additional NOT connected to the output of an AND or an OR gate. The
area will be a little bigger (maybe!), but the strength of the circuit will be
increased because the NOT circuit works as a buffer improving the
time performance of the non-inverting gate.

The propagation time for the 2-input main gates is computed in a similar way as the propagation for NOT circuit is  computed. The only differences are due to the fact that sometimes $R_{ON}$ must be substituted with $2 \times R_{ON}$.

\subsection{Propagation time}

\subsubsection{Propagation time for NAND gate} 

Propagation time becomes, in the worst case when only one input switches:
$$t_{HL} = k_n(2R_{ONn})C_L$$
$$t_{LH} = k_p(R_{ONp})C_L$$
because the capacitor $C_L$ is charged through one pMOS transistor and is discharged through two, serially connected, nMOS transistors.

\subsubsection{Propagation time for NOR gate} 

Propagation time becomes, in the worst case when only one input switches:
$$t_{HL} = k_n(R_{ONn})C_L$$
$$t_{LH} = k_p(2R_{ONp})C_L$$
because the capacitor $C_L$ is charged through two, serially connected, pMOS transistors and is discharged through one nMOS transistor.

\bigskip

It is obvious that we must prefer, when is is possible, the use of NAND gates instead of NOR gates, because, for the same area, $R_{ONp} > R_{ONn}$.

